% Part: lambda-calculus
% Chapter: introduction
% Section: church-rosser

\documentclass[../../../include/open-logic-section]{subfiles}

\begin{document}

\olfileid[fr]{lam}{int}{cr}
\olsection{La propriété de Church--Rosser}

\begin{thm}
\ollabel{thm:church-rosser}
Soient $M$, $N_1$ et $N_2$ des termes tels
que $M\red N_1$ et $M\red N_2$. Il existe
alors un terme~$P$ tel que $N_1\red P$
et $N_2\red P$.
\end{thm}

\begin{cor}
Si $M$ peut se réduire à une forme normale,
cette forme normale est unique à $\alpha$-équivalence près.
\end{cor}

\begin{proof}
Si $M\red N_1$ et $M\red N_2$, le théorème
précédent donne un terme~$P$ auquel $N_1$
et $N_2$ se réduisent tous deux. Si $N_1$
et $N_2$ sont des formes normales, cela
n'est possible que si $N_1\ident P\ident N_2$.
\end{proof}

Enfin, deux termes $M$ et~$N$ sont dits
\emph{$\beta$-équivalents}, ou simplement
\emph{équivalents}, s'ils se réduisent à
un même terme : il existe alors un~$P$
tel que $M\red P$ et $N\red P$. Nous
écrivons $M\equal[\beta] N$. À l'aide du
\olref{thm:church-rosser}, on peut vérifier
que $\equal[\beta]$ est une relation
d'équivalence et que, pour tous $M$ et~$N$,
si $M\red N$ ou $N\red M$, alors
$M\equal[\beta] N$. On peut même montrer
que $\equal[\beta]$ est la \emph{plus petite}
relation d'équivalence possédant cette
propriété.

\end{document}
